\documentclass[11pt]{article}
\usepackage[margin=1in]{geometry}
\usepackage{amsmath,amssymb,amsthm,mathtools}
\usepackage{graphicx}
\usepackage[hidelinks]{hyperref}
\usepackage{microtype}

\newtheorem{theorem}{Theorem}
\newtheorem{lemma}[theorem]{Lemma}
\newtheorem{corollary}[theorem]{Corollary}
\newcommand{\R}{\mathbb R}
\newcommand{\BMO}{\mathrm{BMO}}
\newcommand{\VMO}{\mathrm{VMO}}
\newcommand{\CMO}{\mathrm{CMO}}
\newcommand{\MO}{\operatorname{osc}}
\newcommand{\esssup}{\operatorname*{ess\,sup}}
\newcommand{\fint}{\mathop{\rlap{\raisebox{0.15ex}{\ensuremath{-}}}\!\int}}

\title{Bounded weak solutions have exactly bounded VMO traces}
\author{Same-paper consequence clarifying the VMO bullet of arXiv:2605.00439}
\date{11 August 2026}

\begin{document}
\maketitle

\begin{center}
\fbox{\parbox{0.92\linewidth}{\textbf{Status.}
Verified literal classification implied by the source itself, not claimed as
a new solution of the intended broader VMO program. Under the source's
bounded-solution definition, the admissible VMO traces are
exactly $\VMO\cap L^\infty$; hence no extension to genuinely unbounded VMO
data is possible without changing the solution class.}}
\end{center}

\section{Source problem and answer}

Auscher and Bechtel study
\begin{equation}\label{eq:QL}
 \partial_tu-\operatorname{div}(a(t,x,u)\nabla u)=0,
 \qquad u(0)=u_0,
\end{equation}
on $\R^n$.  Their Definition~1.2 calls $u$ a \emph{bounded weak
solution} only if $u\in L^\infty((0,T)\times\R^n)$ and
\begin{equation}\label{eq:trace}
 \langle u(t),\psi\rangle\longrightarrow\langle u_0,\psi\rangle
 \quad(t\downarrow0),\qquad \psi\in C_c^\infty(\R^n).
\end{equation}
Corollary~1.5 proves existence for every $u_0\in L^\infty$ satisfying
the source range condition.  The paper then asks whether this can be
extended to ``more irregular data such as $u_0\in\VMO(\R^n)$'' \cite{AB}.

\begin{center}
\includegraphics[width=0.98\linewidth]{figures/open_problem_crop.png}
\end{center}

The answer is no for the stated solution class.  In fact there is an exact
classification.

\begin{theorem}[Classification of VMO traces]\label{thm:classification}
Let $a$ and the open set $O$ satisfy the hypotheses of \cite[Assumption
A]{AB}.  Let $u_0\in\VMO(\R^n)$ be a locally bounded representative whose
essential range is contained in $O$.  There exists a bounded weak solution
in the sense of \cite[Definition 1.2]{AB} only if $u_0\in L^\infty(\R^n)$.
Conversely, if $u_0\in L^\infty(\R^n)$ and its closed essential range is
contained in $O$, such a global solution exists.  Thus the source's exact
VMO existence class is
\[
 \VMO(\R^n)\cap L^\infty(\R^n)
\]
(with the same range condition as in the source).
\end{theorem}

The converse is precisely \cite[Corollary 1.5]{AB}.  The necessity is
already contained abstractly in \cite[Proposition 3.4(b)]{AB}; the next
PDE-free trace lemma gives a one-line independent proof.  To show that the
boundary is nonvacuous under every
standard global VMO convention, Section~\ref{sec:example} constructs a
locally bounded but unbounded datum lying even in $\CMO$, the closure of
$C_c^\infty$ in $\BMO$.

\section{The universal bounded-trace obstruction}

\begin{lemma}[Weak-* closedness of the $L^\infty$ ball]\label{lem:trace}
Suppose $u\in L^\infty((0,T)\times\R^n)$ and a distribution $f$ satisfies
\[
 \int_{\R^n}u(t,x)\psi(x)\,dx\longrightarrow\langle f,\psi\rangle
 \quad(t\downarrow0)
\]
for every $\psi\in C_c^\infty(\R^n)$.  Then $f$ is represented by a
function in $L^\infty(\R^n)$ and
\[
 \|f\|_\infty\leq\|u\|_{L^\infty((0,T)\times\R^n)}.
\]
\end{lemma}

\begin{proof}
Put $M=\|u\|_\infty$.  By Fubini, there is a full-measure set
$E\subset(0,T)$ such that $\|u(t,\cdot)\|_\infty\leq M$ for every
$t\in E$.  Choose $t_j\in E$ with $t_j\downarrow0$.  Then
\[
 |\langle f,\psi\rangle|
   =\lim_j\left|\int u(t_j,x)\psi(x)\,dx\right|
   \leq M\|\psi\|_1.
\]
Thus $f$ extends by density to a bounded functional of norm at most $M$ on
$L^1(\R^n)$.  The duality $(L^1)^*=L^\infty$ represents that functional by
some $g\in L^\infty$ with $\|g\|_\infty\leq M$.  Hence $f=g$ as a
distribution.  If the prescribed datum was already a locally integrable
function, it equals $g$ almost everywhere.
\end{proof}

Applying Lemma~\ref{lem:trace} to \eqref{eq:trace} proves the necessary
half of Theorem~\ref{thm:classification}.  Notice that ellipticity,
regularity of $a$, and even equation \eqref{eq:QL} play no role: the
obstruction is built into the simultaneous requirements of uniform
spacetime boundedness and distributional attainment of the initial datum.

\section{A locally bounded unbounded CMO datum}\label{sec:example}

There are two widespread conventions on $\R^n$: Sarason VMO, the
$\BMO$-closure of uniformly continuous BMO functions, and the smaller
$\CMO$, the $\BMO$-closure of $C_c^\infty$.  The following example lies in
the smaller space.

\begin{lemma}\label{lem:cmo}
For every $n\geq1$, with $A=e^e$, the radial function
\begin{equation}\label{eq:f}
 f(x)=\log\log(A+|x|)
\end{equation}
is locally bounded, unbounded, and belongs to $\CMO(\R^n)$.
\end{lemma}

\begin{proof}
The first two claims are immediate.  We verify Uchiyama's three
mean-oscillation conditions, which characterize the $\BMO$-closure of
$C_c^\infty$ \cite{Uchiyama}.  Balls and cubes give equivalent conditions.
For a ball $B$, write
\[
 \MO_B(f)=\frac1{|B|}\int_B|f-f_B|.
\]

First, $f$ is globally Lipschitz: away from the origin,
\[
 |\nabla f(x)|=
 \frac1{(A+|x|)\log(A+|x|)}\leq\frac1{Ae},
\]
and the same bound follows across the origin from the radial profile.
Consequently
\begin{equation}\label{eq:small}
 \sup_{\operatorname{rad}(B)\leq r}\MO_B(f)\leq\frac{2r}{Ae}\to0
 \qquad(r\downarrow0).
\end{equation}

We next estimate centered balls.  Put $h(r)=\log\log(A+r)$ and
$L_R=\log(A+R)$.  In radial coordinates $s=|x|/R=e^{-y}$,
the normalized measure is $n e^{-ny}\,dy$.  Since
$A+Re^{-y}\geq e^{-y}(A+R)$, for $0\leq y\leq L_R/2$,
\[
 0\leq h(R)-h(Re^{-y})
 \leq-\log(1-y/L_R)\leq\frac{2y}{L_R}.
\]
For $y>L_R/2$, use the crude bound
$h(R)-h(Re^{-y})\leq\log(L_R/e)$.  It follows that
\begin{equation}\label{eq:centered}
 \fint_{B(0,R)}|f-h(R)|
 \leq \frac{2}{nL_R}+\log(L_R/e)e^{-nL_R/2}\longrightarrow0.
\end{equation}
Mean oscillation is at most twice the average distance from any constant.

Now let $B=B(c,r)$ with $r\to\infty$.  If $|c|\leq2r$, then
$B\subset B(0,3r)$ and comparison with the constant $h(3r)$, followed by
\eqref{eq:centered}, gives
\[
 \MO_B(f)\leq 2\,3^n\fint_{B(0,3r)}|f-h(3r)|\longrightarrow0.
\]
If $|c|>2r$, the gradient bound on $B$ instead gives
\[
 \MO_B(f)\leq
 \frac{4r}{|c|\log(A+|c|/2)}
 \leq\frac{2}{\log(A+r)}\longrightarrow0.
\]
This proves uniform vanishing on large balls.  Finally, for any fixed ball
translated to infinity, the same gradient estimate makes its oscillation
tend to zero.  Together with \eqref{eq:small}, these are Uchiyama's three
conditions.  They also imply $f\in\BMO$ (the remaining bounded range of
radii and centers is compact, while the far-center gradient bound is
uniform).  Therefore $f\in\CMO(\R^n)$.
\end{proof}

\begin{corollary}[Explicit counterexample]\label{cor:counterexample}
Take $O=\R$ and $a(t,x,y)=I_n$.  Then \eqref{eq:QL} is the heat equation
and satisfies every source assumption.  For the locally bounded datum
\eqref{eq:f}, there is no bounded weak solution in the source sense.
\end{corollary}

\begin{proof}
Lemma~\ref{lem:cmo} gives $f\in\CMO\subset\VMO$ and
$f\notin L^\infty$.  Lemma~\ref{lem:trace} rules out the required bounded
solution before the heat equation is used.
\end{proof}

\section{Scope and interpretation}

The result does not say that VMO data cannot seed a useful parabolic
theory.  It says that genuinely unbounded VMO data cannot do so while the
solution remains essentially bounded on the whole cylinder down to
$t=0$ and attains the datum distributionally.  A meaningful extension
must therefore change at least one ingredient, for example by allowing a
time-dependent bound that diverges at the initial time or by using a
BMO/VMO-adapted solution norm.  If ``VMO data'' was intended to mean only
$\VMO\cap L^\infty$, existence is already contained in Corollary~1.5 of
the source.

The current v3 source already proves that every bounded linear solution has
an $L^\infty$ trace in Proposition~3.4(b), and every bounded quasilinear
solution is such a linear solution after freezing its coefficient. Thus this
note is a same-paper logical consequence, not a novelty claim. Its purpose is
to identify precisely how any nontrivial VMO reformulation must differ.

\begin{thebibliography}{9}
\bibitem{AB}
P.~Auscher and S.~Bechtel,
\emph{Existence and uniqueness of weak solutions to quasilinear PDEs with
critical data}, arXiv:2605.00439v3, 2026.

\bibitem{Uchiyama}
A.~Uchiyama,
\emph{On the compactness of operators of Hankel type},
Tohoku Math. J. (2) \textbf{30} (1978), 163--171; see the CMO
characterization on p.~166.
\end{thebibliography}

\end{document}
